% Generated by report_gen.py (ai-ee v1) - do not hand-edit.
\documentclass[11pt,a4paper]{article}
\usepackage[margin=2.2cm]{geometry}
\usepackage{graphicx}
\usepackage{booktabs}
\usepackage{longtable}
\usepackage{pdfpages}
\makeatletter
\@ifundefined{KV@Gin@artifact}{\define@key{Gin}{artifact}[]{}}{}
\makeatother
\usepackage{xcolor}
\usepackage[hidelinks]{hyperref}
\setcounter{tocdepth}{1}
\begin{document}
\sloppy
\begin{center}
{\LARGE\bfseries bb-buck --- Design Document}\\[6pt]
{\large ai-ee v1 pipeline}\\[2pt]
generated 2026-08-15 22:22:41
\end{center}
\tableofcontents

\section{Board and Run Metadata}
{\small
\begin{longtable}{lp{11cm}}
\toprule
\textbf{Field} & \textbf{Value} \\
\midrule
\endhead
board & bb-buck \\
workspace & \texttt{boards/bb-buck} \\
phase & P4 \\
gate status & no gates recorded yet \\
state created & 2026-08-15T15:58:13 \\
state updated & 2026-08-15T22:11:49 \\
\bottomrule
\end{longtable}
}
\section{Overview}
\subsection*{Brief - bb-buck (as given by the owner, verbatim)}
ultra bare bones design: a switch-mode step-down converter. Input 24 V DC nominal (18-30 V) on a screw terminal; output 5 V at up to 2 A on a screw terminal. Bench supply in, resistive load out. Workspace boards/bb-buck.

\subsubsection*{Orchestrator notes (not owner words)}
\begin{itemize}
\item Mode token present: \texttt{ultra bare bones design:} -\textgreater{} the ultra-bare-bones scope
contract in \texttt{reference/build-modes.md} applies.
\item Block under study: the synchronous/asynchronous step-down (buck) converter.
\item Bench-only article: bench DC supply on the input, resistive load on the output.
\end{itemize}
\subsection*{Stackup}
Chosen stackup: \textbf{`JLC2313\_1.6`} - JLCPCB standard 2-layer, 1.6 mm, \textbf{1 oz outer copper},
\section{Requirements}
\subsection*{Requirements: bb-buck}
\subsubsection*{1. Function}
A bare step-down (buck) DC-DC converter, built as a study article. It takes 18-30 V DC (24 V nominal) from a bench power supply on a 2-pin screw terminal and produces 5 V at up to 2 A (10 W) on a second 2-pin screw terminal feeding a resistive load. Non-isolated, common ground in to out. Nothing else is on the board: no MCU, no second rail, no digital interface, no indicators.

BUILD MODE: \textbf{ultra-bare-bones} (declared by the brief's opening token, recorded as a P0 decision; contract in \texttt{reference/build-modes.md}). The block under study is the buck converter itself. Scope is that block plus exactly the support parts its datasheet requires at the stated operating point, one input and one output interface, and what the fab needs to build the board and the bench needs to hold it. Protection, filtering, indicators, spare rails and enclosure features are OUT by mode, not by engineering judgment. The mode's defaults are applied below (quantity 5, JLC PCBA single-sided, bench 0-50 C indoor, no enclosure, fewest honest layers, smallest honest outline, stop at P9) instead of being asked.

Synchronous vs asynchronous rectification is deliberately NOT fixed here. The brief leaves it open and it is a P1 architecture call, not a requirement. Both are in scope as long as the block meets the stated operating point.

\subsubsection*{2. Interfaces}
Exactly two, both stated by the brief.

\begin{itemize}
\item Power input (J1): 2-pin screw terminal (stated). DC only, 18-30 V.
ASSUMED: through-hole 2-pole terminal block, 5.08 mm pitch, rated
\textgreater{}= 300 V and \textgreater{}= 10 A so it is never the limiting element against the
\textasciitilde{}0.7 A worst-case input current. Pitch was not stated; 5.08 mm is the
common, well-stocked JLC/LCSC choice and accepts bench-supply lead wire
comfortably. Polarity silkscreen-marked (+ / -). NOTE: the board has no
reverse-polarity protection by mode, so the silk marking is the only
defense against a swapped supply - see section 8.
\item Power output (J2): 2-pin screw terminal (stated). 5 V, up to 2 A.
ASSUMED: same terminal family and part as J1 for BOM consolidation,
placed on a different board edge and silk-labelled distinctly (VIN vs
5V OUT) so the two cannot be confused at the bench.
\item No other external interface exists or is wanted: no USB, no header, no
comms bus, no external enable, no power-good, no fault output, no LED.
All excluded by mode.
\item Measurement access ANSWERED (A4): exactly one switch-node probe pad with
an adjacent ground pad. No other test points - VIN, VOUT and GND are
reachable at the terminals.
\end{itemize}
\subsubsection*{3. Power}
\begin{itemize}
\item Source: bench DC power supply, 18-30 V, 24 V nominal (STATED
explicitly by the brief - not a guess, not a battery, not mains-derived
on this board, not automotive). Section 8 depends on this being stated.
\item Load: resistive, up to 2 A at 5 V (STATED). ASSUMED: the load range is
0 to 2 A, i.e. the board must also be stable and safe with NO load
connected (the realistic bench case of powering up before the resistor
is attached). This is the conservative reading of "up to 2 A" and it
constrains the converter's light-load behavior, so it is recorded rather
than left silent.
\item ASSUMED: 2 A is the absolute maximum output, continuous. A resistive
load draws no surge beyond its steady value, so no peak allowance above
2 A is carried. Inrush into the board's own output capacitance at
startup is the converter's soft-start problem, not a load requirement.
\item Output rail: one rail, fixed 5 V nominal. ASSUMED: fixed output, not
adjustable, no trim, no sequencing. Accuracy and ripple ANSWERED (A3):
+/-3 \% DC (4.85-5.15 V) across the full 18-30 V input and 0-2 A load
range, ripple \textless{}= 50 mV peak-to-peak.
\item Input current (GUESS, for sizing only): 10 W out at an assumed 88-93 \%
efficiency -\textgreater{} 10.8-11.4 W in. At 18 V low line \textasciitilde{}0.60-0.63 A; at 24 V
\textasciitilde{}0.45-0.47 A; at 30 V high line \textasciitilde{}0.36-0.38 A. Size the input terminal,
input copper and input capacitor RMS rating for \textgreater{}= 1 A continuous with
a switched peak of \textasciitilde{}2.5 A. Output terminal and copper for \textgreater{}= 2.5 A
continuous.
\item Board dissipation (GUESS): \textasciitilde{}0.8-1.4 W typical, \textasciitilde{}1.8 W at a pessimistic
85 \% efficiency. Worst case sits at the 30 V high-line corner where
switching loss dominates. At the mode's 50 C maximum ambient with
natural convection this is a real thermal driver on a small outline -
copper area under the converter is a design input, not an afterthought.
\item Duty cycle (derived, for P1): D \textasciitilde{} Vout/Vin = 0.28 at 18 V, 0.21 at
24 V, 0.17 at 30 V. The 30 V corner sets the minimum on-time
constraint, which bounds how high the switching frequency can go for a
given part - a real P1/P2 selection constraint, flagged here, not
decided here.
\item No battery, no charging, no energy storage function, no backup rail.
None stated and none implied by a bench supply source.
\item Excluded by mode (SCOPE decisions, not engineering conclusions - do not
record any of these as generally unnecessary): input TVS/clamp, reverse
polarity protection, input fuse, pi/EMI input filter, OVP/OCP/UVLO
beyond what lives inside the converter IC, output LC filter beyond what
the datasheet requires.
\end{itemize}
\subsubsection*{4. Environment}
Mode defaults - applied, not asked.

\begin{itemize}
\item Ambient: 0 to 50 C, indoor bench.
\item Cooling: natural convection only. No forced airflow, no heatsink, no
chassis thermal path.
\item No enclosure. No ingress (IP) rating. No vibration or shock
requirement. No humidity requirement. No conformal coating.
\item No formal EMC/emissions campaign - this is a bench study article, not a
product. Tight power-loop layout is still required because it is what
makes the block correct, not because a test demands it.
\end{itemize}
\subsubsection*{5. Size \& mounting}
\begin{itemize}
\item Outline: \textbf{no HARD cap.} Nothing here binds permanently at P5
board\_init. The mode governs instead: the smallest outline that keeps
the layout HONEST - never so tight that the power loop, the copper
thermal area or the terminal spacing stop being representative of a
real 24 V -\textgreater{} 5 V / 2 A buck, and never padded for features this board
does not have.
\item Non-binding expectation for planning only (NOT a cap): roughly
30-40 mm x 20-28 mm, dominated by the two screw terminals plus the
converter, inductor and capacitors. If the honest layout needs more,
it takes more.
\item Layers: ASSUMED 2, per the mode's "fewest layers the block honestly
needs". If the P2/P4 thermal work shows 2 layers cannot carry \textasciitilde{}1.2-1.8 W
at 50 C ambient inside the honest outline, 4 layers is allowed under the
same rule. Not fixed here either way.
\item Mounting: ASSUMED 4 x M3 clearance holes (3.2 mm) inset from the
corners, so the board can sit on standoffs at the bench. Mounting the
bare board is explicitly in mode scope. If the honest outline turns out
too small for four, two on opposite corners is acceptable.
\item Height: ASSUMED no maximum. The board is open on a bench with no
enclosure, and the screw terminals themselves stand \textasciitilde{}10-12 mm tall with
wire fitted, so a height cap would be arbitrary.
\item Connector placement: ASSUMED both terminals on board edges with their
wire openings facing outward, on different edges, so bench wiring does
not cross the board. Which edges is the layout engineer's call.
\end{itemize}
\subsubsection*{6. Quantity \& budget}
Mode defaults - applied, not asked.

\begin{itemize}
\item Build quantity: 5.
\item Cost: minimal, with no hard per-unit cap. Prefer JLC Basic/Preferred
parts wherever a Basic part actually meets the electrical requirement;
Extended is allowed for the converter IC and the power inductor, where
Basic stock is thin in this class.
\end{itemize}
\subsubsection*{7. Assembly}
Mode defaults - applied, not asked.

\begin{itemize}
\item JLCPCB PCBA, single-sided: all SMT parts on the top side only. The
bottom side stays clear of SMT so it can serve as unbroken
thermal/return copper.
\item The two screw terminals are through-hole. ASSUMED: fitted by JLCPCB
through-hole assembly if offered at acceptable cost for the chosen
part, hand-soldered after SMT otherwise. Either path leaves the
schematic and the footprint unchanged.
\item Pipeline stops at P9 (fab package + DFM). Ordering is a separate owner
decision and is not requested here.
\end{itemize}
\subsubsection*{8. Compliance/safety flags}
No safety question is open on this board, and that conclusion is derived, not assumed. The mode grants no silence here - it simply happens that the brief states the two facts that would otherwise have to be asked.

\begin{itemize}
\item \textbf{Mains voltage: DOES NOT APPLY.} The board sees DC only. Generating
that DC is out of scope; the brief names a bench supply as the source.
\item \textbf{Batteries: DOES NOT APPLY.} The source is STATED as a bench power
supply. This is the one question that has no default and must never be
guessed, and the brief answers it explicitly, so it is not asked. No
battery on the board, no charging function, no pack chemistry, no
deep-discharge or UVLO behavior required. If the owner ever runs this
board from a battery instead, that is a new requirement and this
section must be re-opened.
\item \textbf{Motors: DOES NOT APPLY.} The load is STATED as resistive. No
inductive, stalling or regenerative load, therefore no reverse energy
into the output and no startup surge beyond the resistor's own draw.
\item \textbf{RF transmit: DOES NOT APPLY.} No radio function.
\item \textbf{High current (\textgreater{}3 A): DOES NOT APPLY.} Maximum continuous current
anywhere on the board is the 2 A output. Input current peaks at \textasciitilde{}0.63 A
continuous at the 18 V low-line corner. The largest instantaneous
current is the inductor peak, GUESSED at \textasciitilde{}2.3-2.6 A depending on the
ripple ratio P1 chooses - still under 3 A. The threshold is not
approached closely enough to be a judgment call.
\item {*}{*}\textgreater{}30 V: DOES NOT APPLY at the stated number, but the range's top edge
is genuinely under-defined.{*}{*} Reasoning, explicitly:
\begin{itemize}
\item The stated maximum is 30 V, and the threshold is \emph{greater than} 30 V.
30 is not \textgreater{}30, so the flag does not trip on the stated figure. This
is a boundary case, so it is reasoned about rather than waved past.
\item No real hazard exists at 30 V DC. It is well below the 60 V DC limit
for a non-hazardous DC source under IEC 62368-1 (ES1), so there is no
shock hazard to a person handling the board. Fault energy is bounded
by the bench supply's own current limit, and at \textless{}= 1 A input there is
no arc or energy hazard at the terminals. No safety question follows
from the voltage itself.
\item Downstream consequence worth noting: at exactly 30 V, the VIN-to-GND
net pair is 30 V apart, which the pipeline's \texttt{check\_creepage} treats
as a clean no-op (its derived check engages only above 30 V). If the
answer to Open question 1 lifts the ceiling above 30 V, that check
engages for real and IPC-2221 spacing starts binding the layout. The
boundary is therefore not cosmetic.
\item What WAS genuinely ambiguous - whether 30 V is the \emph{operating} or the
\emph{survival} ceiling, and what the input does on a live connection - is
now ANSWERED (A1, A2): 30 V is a hard maximum OPERATING input, and
there is no live hot-plug. The converter carries \textgreater{}= 36 V absolute-max
as its own headroom; no clamp is added. The flag therefore stays
DOES-NOT-APPLY on settled requirements, not merely on the boundary
reading.
\end{itemize}
\item \textbf{Reverse polarity: noted, not flagged.} The mode excludes reverse
polarity protection, so connecting the bench supply backwards will
destroy the converter. This is an accepted consequence of the declared
scope on a bench article with silk-marked polarity, not a safety hazard
to a person, and not a finding for a reviewer to raise. It is recorded
so the consequence is visible rather than implicit.
\end{itemize}
\subsubsection*{9. Open questions - ALL FOUR ANSWERED 2026-08-15 (owner)}
STATUS: \textbf{closed.} The owner took the stated DEFAULT on all four. Each answer is recorded as a P0 \texttt{state.py decision} and is now a REQUIREMENT, not a question. P1/P2 read the ANSWERED block below; the original question text is kept underneath it for provenance only.

\paragraph*{Answers (binding)}
\begin{itemize}
\item \textbf{A1 - Input ceiling.} 30 V is a HARD maximum operating input. Choose a
converter whose absolute-maximum input rating is \textgreater{}= 36 V class (\textasciitilde{}20 \%
headroom above 30 V). The headroom lives in the PART, not in added
components - no clamp, no TVS. Consequence retained: VIN-to-GND stays at
or below 30 V, so \texttt{check\_creepage} remains a no-op and IPC-2221 spacing
does not bind the layout.
\item \textbf{A2 - No live hot-plug.} The input is never connected while the supply
is live: wires land first, then the supply ramps from zero. The \textasciitilde{}2x
lead-inductance ringing case (\textasciitilde{}60 V from a 30 V setting) is therefore OUT
of the requirement set. The datasheet-required bulk input capacitance
stays fitted and damps the incidental case; no damping component is
added. A 36 V-class part is sufficient under this answer.
\item {*}{*}A3 - Output spec. AMENDED TWICE AT P3/P4 - read the amendment below, not
just this paragraph.{*}{*} As originally answered: +/-3 \% DC, i.e. 4.85-5.15 V,
held across the FULL 18-30 V input range and the FULL 0-2 A load range,
with \textless{}= 50 mV peak-to-peak output ripple.
\textbf{A3 AS AMENDED (this is the binding version).} A3 now has a TWO-REGION
shape, and both amendments have the same root cause: the chosen part runs
PFM / diode-emulation at light load and has no MODE pin to defeat it.
\end{itemize}
\begin{quote}\small\ttfamily\raggedright
  \textbar{} Load \textbar{} DC accuracy \textbar{} Ripple \textbar{}\\
  \textbar{}---\textbar{}---\textbar{}---\textbar{}\\
  \textbar{} \textasciitilde{}200 mA to 2 A \textbar{} +/-3 \% (4.85-5.15 V) BINDING \textbar{} \textless{}= 50 mV pk-pk BINDING \textbar{}\\
  \textbar{} 0 to \textasciitilde{}200 mA \textbar{} PFM light-load rise ACCEPTED \textbar{} PFM burst ripple ACCEPTED \textbar{}
\end{quote}
\begin{itemize}
\item \emph{Amendment 1 (P3, owner-ruled).} The DDA package has no MODE/FCCM pin and
is auto-mode only (datasheet 8.4.1, which states outright that PFM
"yields larger output voltage ripple"). None of the three accept criteria
\begin{itemize}
\item forced-CCM, a strappable MODE pin, or a datasheet figure at this
operating point - was satisfiable. The mode excludes the usual fix (a
preload resistor is excluded conditioning), so the owner relaxed ripple
below \textasciitilde{}200 mA rather than override the mode or re-source the part.
\item \emph{Amendment 2 (P4, owner-ruled).} The SAME PFM behaviour also raises the
output VOLTAGE at light load: datasheet Sec 7.7 specs regulation as
-1.5/+1.5 \% for IOUT \textgreater{}= 1 A but \textbf{-1.5/+2.5 \% for IOUT 0 A to max}, a
+1.0 \% adder confirmed by Fig 9-19. Stacked on the worst 0-50 C divider
corner that gave 5.161 V against the 5.15 V ceiling - outside. So the
carve-out extends to DC accuracy below \textasciitilde{}200 mA. An earlier statement that
DC accuracy was unaffected at all loads was WRONG and is retracted here.
\item \emph{Paid for, not merely relaxed.} The owner took the recentring remedy as
well: the feedback divider moved from 100k/24.9k (5.016 V nominal, TI's
own Table 9-2 pair) to \textbf{102k/25.5k}, an exact 4.000 ratio giving
\textbf{5.000 V nominal}. That moved the stacked 0 A worst case from 5.161 V
to 5.144 V - inside the original window even though it no longer has to
be. Both parts are E96, 0.1 \%/25 ppm, same Yageo RT family (TCR tracking
preserved).
The P8 simulation bounds encode this: error bounds stay on the full
4.85-5.15 V window (the divider ratio is load-independent, so the carve-out
relaxes the CONVERTER, not the divider), plus a warning-severity bound
carrying the +1.0 \% light-load stack so it stays visible in the gate. That
warning bound is the ONLY check that catches a revert to 100k/24.9k.
\end{itemize}
\item \textbf{A4 - Probe access.} Exactly ONE switch-node probe pad with an adjacent
ground pad. Nothing else: VIN, VOUT and GND are already reachable at the
two screw terminals. Mode-boundary ruling by the owner - the switch node
counts as the block's own measurement need on a study article. Keep the
pad minimal and inside the switch-node copper that already exists, so it
adds as little area as possible to the noisiest node.
\end{itemize}
\paragraph*{Original question text (provenance)}
1. {*}{*}Is 30 V truly the highest the supply will ever be set to, or should the board survive a bit more?{*}{*} A bench supply set to "30 V" may actually sit at 30.5 V, and a knob can be nudged. This board has no input clamp by design, so the converter chip's own maximum rating is the only thing standing between an overshoot and a dead board. DEFAULT: 30 V is a hard maximum operating input, and the converter is chosen with an absolute-maximum input rating at least \textasciitilde{}20 \% above it (36 V class or better) purely as selection headroom. No extra parts are added. 2. {*}{*}Will you ever connect the input wires while the bench supply is already switched on and at voltage?{*}{*} Connecting live through a metre of lead wire makes the input voltage ring up to roughly twice the supply setting for a few microseconds - about 60 V from a 30 V supply - because the lead inductance rings against the board's ceramic input capacitors. That transient is well beyond what a 36-40 V part tolerates, and the mode excludes the clamp that would normally absorb it, so the answer changes which parts are even eligible. DEFAULT: no live hot-plug. You connect the wires first, then bring the supply up from zero. The bulk input capacitance the datasheet already requires stays fitted and damps the incidental case; nothing extra is added. If you answer "yes, I will hot-plug it", say so - it will force either a much higher-voltage part or a damping component the mode would otherwise exclude, and that is an owner call. 3. \textbf{How tight does the 5 V output need to be?} Two numbers: how far from exactly 5.00 V is acceptable, and how much high-frequency ripple can ride on top of it. DEFAULT: +/-3 \% DC (4.85-5.15 V) across the full 18-30 V input range and the full 0-2 A load range, and \textless{}= 50 mV peak-to-peak output ripple. Ordinary, honest numbers for a study article - not a precision rail, not a loose one. 4. \textbf{Do you want one probe point on the switching node?} That means a small pad on the switch node with a ground pad right beside it, so a scope probe with a short ground can measure switch-node ringing and efficiency properly at bring-up. This is the one place the mode's boundary needs an owner ruling: the mode excludes test points "beyond the block's own measurement need", and for a board whose purpose is to be studied - and whose bring-up evidence is what proves out the knowledge records it exists to seed - that need is arguable either way. The tradeoff is real: a pad hangs extra copper on the noisiest node on the board. DEFAULT: yes - exactly one switch-node probe pad with an adjacent ground pad, and nothing else. VIN, VOUT and GND are already reachable at the two screw terminals, so no other test points are added.

\section{Architecture}
\subsection*{Decisions of record (state.json)}
{\small
\begin{longtable}{lp{5.2cm}p{7.2cm}}
\toprule
\textbf{Phase} & \textbf{What} & \textbf{Why} \\
\midrule
\endhead
P0 & Build mode: ultra-bare-bones (token 'ultra bare bones design:' opens the brief) & reference/build-modes.md scope contract: ONE functional block (24V-\textgreater{}5V/2A buck) plus only datasheet-required support; one input + one output interface (screw terminals); no protection/filtering/indicators/test points beyond the block's own need; defaults = fewest honest layers, JLC PCBA single-sided, qty 5, bench 0-50C, stop at P9. Research mandatory: first coverage call of each phase runs --maturity-floor proven --research-provisional. \\
P0 & Input ceiling: 30 V is a HARD maximum operating input; converter chosen with Vin absolute-max \textgreater{}= 36 V class (\textasciitilde{}20\% headroom) as selection headroom only - no added clamp parts & Owner answered P0 Q1 with the default. Keeps VIN-GND at/below 30 V so check\_creepage stays a no-op and IPC-2221 spacing does not bind the layout; the mode excludes an input clamp so the IC rating is the only defense. \\
P0 & No live hot-plug of the input: wires land first, supply ramps from zero & Owner answered P0 Q2 with the default. Removes the \textasciitilde{}2x lead-inductance ringing case (\textasciitilde{}60 V from 30 V) from the requirement set; datasheet-required bulk input capacitance damps the incidental case, no damping component added. \\
P0 & Output spec: +/-3\% DC (4.85-5.15 V) over 18-30 Vin and 0-2 A load; ripple \textless{}= 50 mV peak-to-peak & Owner answered P0 Q3 with the default. Sets the regulation/ripple targets that bind feedback divider tolerance, output capacitance and the P8 sim bounds. \\
P0 & One switch-node probe pad plus an adjacent GND pad; no other test points & Owner answered P0 Q4 with the default. Mode-boundary ruling: the switch node is the block's own measurement need on a study board whose bring-up evidence proves the seeded records. Kept minimal and inside the existing SW polygon to limit added copper on the noisiest node. \\
P1 & P1 research roster: research-power-architect + research-component-scout(buck-regulator) only. NO research-interface-spec, NO research-reference-design. & Screw terminals are bound by no electrical standard, so there is no interface spec to research. The buck block is not novel - the library already holds 16 buck records - and the useful vendor-layout digest is exactly what the P3 part-level coverage research produces for the CHOSEN IC; a P1 reference-design spawn would duplicate it. Mode says do not pad. \\
P1 & Topology: synchronous integrated-FET buck, Fsw \textasciitilde{}500 kHz (band 350-600), L = 10 uH & research-power-architect numbers: async costs +0.42 W (+33\%) and adds a 0.77 W diode hot spot at Tj \textasciitilde{}126 C against a 125 C rating, worst-cased by D=0.17 (diode conducts 83\% of every cycle at 30 V). Fsw \textgreater{}= 1 MHz costs 1.58-1.93 W and collides t\_on(30 V) with typical min-on-time. LDO is 38 W - arithmetic. Controller+FETs is the ladder's escape hatch for corners with no stocked integrated part; 30 V/2 A is not one. \\
P1 & PROVISIONAL 2-layer stackup, conditional on the P3 part: sync + Rds\_LS \textless{}= 85 mOhm + outline \textgreater{}= \textasciitilde{}1000 mm2 + \textgreater{}= 16 thermal vias + Tj 150 C & The 2L screen passes by only 2 C (0.92 W x 73.8 C/W = 67.9 vs check\_thermal dt\_c 70). The layer count FOLLOWS THE PART - a 2 A-class part at Rds\_LS 130 mOhm gives 1.19 W and forces 4L. P3 must re-run this arithmetic with the chosen part's real Rds before P5 fixes the stackup (no outline-shrink step exists; board\_init is final). \\
P1 & OPEN CONSTRAINT carried to P3 (not settled at P1): light-load ripple vs A3. A3 binds \textless{}= 50 mVpp across the FULL 0-2 A range INCLUDING no load, but PFM/burst light-load ripple is typically 50-100 mVpp at 5 V. & research-power-architect OPEN item. The mode excludes the usual fix (a preload resistor is excluded conditioning), so P3 part selection must supply ONE of: a forced-CCM part, a strappable MODE/FCCM pin, or a datasheet figure showing \textless{}= 50 mVpp at light load. If no JLC-stocked 36 V-class candidate offers any of the three, ESCALATE to the owner to choose between relaxing A3 at light load and admitting a preload against the mode. A burst-mode part must NOT ship silently against A3. \\
P2 & SKIPPED the optional coverage-mapper spawn on the P2 mapping\_request & Every one of the 15 buck records is ALREADY matched to block:B1 via keys and evaluated per class; the blockers are maturity-below-floor (8 classes), envelope-outside (inrush on source\_kind, snubber on rectifier\_kind, COT on control\_kind) and level-below-min (1). A mapping edge cannot change any of those three verdicts, so a mapper pass is a provable no-op here. The only unkeyed library record (spiral-inductor-plane-and-heatsink-keepout) is also 'approved' and would still be maturity-blocked. Recipe step marks the mapper optional. \\
P2 & D5 OUTLINE 40 x 30 mm (1200 mm2) - FINAL, binds at P5 board\_init & The outline IS the radiator: R\_ba is 39/34/31 C/W at 875/1064/1200 mm2, so area buys more Tj than layers (\textasciitilde{}10 C vs \textasciitilde{}11 C) and is free at the fab. 20\% above P1's \textasciitilde{}1000 mm2 floor. The increment over 38x28 is bought by physical parts - a 10-12.5 mm inductor (D11), two \textasciitilde{}10 mm screw terminals, 4x M3 with 6.5 mm keepouts - not padding. No outline-shrink step exists later. J1 LEFT edge, J2 BOTTOM edge. \\
P2 & D4 STACKUP JLC2313\_1.6, 2 layers, 1 oz - CONDITIONAL & Conditions: sync + P\_U1 \textless{}= 0.95 W at 30V/2A + outline \textgreater{}= 1000 mm2 with unbroken B.Cu pour + \textgreater{}= 16 thermal vias + Tj\_max 150 C. Named escalation trigger: P\_U1 \textgreater{} 0.95 W (at 400 kHz = Rds\_LS \textgreater{} \textasciitilde{}110 mOhm at 25 C; at 500 kHz, 85 mOhm) OR a 125 C part OR an async part OR check\_thermal erroring at dt\_c 70 -\textgreater{} 4L JLC04161H-1080B, outline unchanged. P3 MUST re-run this with the chosen part's real datasheet Rds BEFORE P5 freezes the stackup. \\
P2 & D2 Lead part class: 36 V-class SYNC integrated-FET 3 A-class, exposed-pad SOIC-8-EP, fixed 400 kHz, internally compensated, peak-current-mode (LMR33630ADDAR). ADJUSTABLE output accepted. & CONFLICT RESOLVED: P1 preferred a FIXED 5 V part to delete the divider tolerance term; the scout's JLC sweep found no fixed 5 V sync 36 V-class part, and its only fixed candidate (XL1509-5.0E1) is +/-4 \% guaranteed, which alone breaks A3. P1's PREFERENCE loses to stock reality - it was a preference, not a filter. Scout ranks 2-4 are asynchronous and are NOT drop-in fallbacks (D1 rejects async with numbers). If P3 finds a stocked fixed 5 V sync \textgreater{}= 36 V part, take it: R1/R2, the tolerance budget and the only sim candidate all disappear. \\
P2 & D3 fsw 400 kHz, L = 15 uH (P1 recommended 500 kHz / 10 uH) & The lead part is a FIXED 400 kHz device with no RT pin, and 400 kHz sits inside P1's own 350-600 band. At 400 kHz L must be 15 uH (10 uH gives 52\% ripple); P\_U1 improves 0.924 -\textgreater{} 0.854 W. dI 0.694 A, I\_L,pk 2.35 A, ripple 10.4 mVpp. Reverses if the part's internal-compensation L/C table excludes 15 uH - fallback 10 uH at 400 kHz (15.6 mVpp, I\_L,pk 2.52 A, both in spec). \\
P2 & D8 canonical net names +VIN /SW +5V GND (/FB and /BST deliberately NOT in constraints power{[}{]}) & P1 proposed VIN/SW; renamed to match the KiCad export mechanism (power symbols export bare, root-sheet local labels export with one leading slash). A power{[}{]} entry on /FB would put a width rule on a high-impedance sense node and pull it into the check\_pdn decoupling inventory. P4 must run netlist\_audit.py - a rename surfaces as missing\_net (error) and nothing else catches it. \\
P2 & D9 ONE FLAT ROOT SHEET, hierarchy rejected; D6 FB divider 0.1 \%/25 ppm; D12 exactly ONE sim candidate (FB divider DC accuracy vs A3's window) & D9: 14 parts / 6 nets, and a child-sheet label exports as /\textless{}sheet\textgreater{}/\textless{}LABEL\textgreater{}, silently unhooking every constraints entry spelling it /\textless{}LABEL\textgreater{} (the class that cost lumina-par). D6: 1 \% resistors give a worst-case SUM of \textasciitilde{}3.7 \% which fails A3; 0.5 \% closes only on RSS; 0.1 \%/25 ppm closes on the worst-case SUM (\textasciitilde{}2.5 \%) for cents and no part-count change. D12: buck switching is not simulated by policy; the FB divider sweep catches the wrong-value class ERC/DRC/DFM are blind to, and disappears if P3 lands a fixed-output part. \\
P3 & P3/P6/P7 spawns receive the knowledge prompt\_block BY PATH (log/knowledge-P3.json, field prompt\_block) instead of pasted inline & knowledge.py --select returned 24 records / 37067 chars. Pasting that inline into every P3/P6/P7 spawn would consume the orchestrator's context for content the orchestrator must never reason about anyway (rule 1: agents read the design, the orchestrator reads contracts). Files are the interface - the agent reads the same bytes either way, and the selection is sha-stable on disk. Deviation from the SKILL.md 'paste its prompt\_block here' wording, recorded. \\
P3 & D4 CONDITION MET - 2 LAYERS CONFIRMED. JLC2313\_1.6 stands; no escalation to 4L. & Real datasheet Rds (LMR33630 DDA, p7 Electrical Characteristics): Rds\_LS 66 mOhm typ / 110 mOhm MAX guaranteed across -40..+125 C (already worst-case, no extra hot-derating owed); Rds\_HS 95/160. Conduction D{*}I\textasciicircum{}2{*}Rds\_HS + (1-D){*}I\textasciicircum{}2{*}Rds\_LS at I\_L,rms 2.01 A, D 0.167 = 286 mW typ / 478 mW worst-case; plus power\_tree's switching+gate+Iq 286 mW = P\_U1 0.57 W typ / 0.76 W worst-case against the 0.95 W gate. \textasciitilde{}20\% margin at worst case, versus the 2-7 C cushion the placeholder model predicted. The stackup is now free to freeze at P5. \\
P3 & D3 CONDITION MET - L = 15 uH and a 2 x 22 uF C\_OUT bank both sit inside the part's windows. No fallback to 10 uH needed. & 15 uH clears the 3.5 uH slope-comp floor and sits at 23\% ripple against the device max (inside the 20-40\% band). Correction to the knowledge record's reading: the 52/72 uF C\_OUT figure is TI's own LOAD-TRANSIENT sizing choice (datasheet 9.2.2.5 - 'usually limited by load transient requirements, rather than output voltage ripple'), not a stability floor. bb-buck carries no binding transient spec (power\_tree s5 already calls the resulting sag normal), so the binding constraint is the ripple equation, already at 10.4 mV. \\
P3 & D11 CONDITION MET - inductor Isat 8 A vs the 6.6 A floor, DCR 27 mOhm, body 12.3 x 12.3 x 8.0 mm & Floor is 1.3 x the part's max peak current limit I\_SC 5.05 A = 6.6 A, not 1.3 x the 2.35 A operating peak. 21\% margin. DCR 27 mOhm beats the 30 mOhm preferred target. Body sits at the TOP of the 10-12.5 mm class the 40 x 30 mm outline was sized around - it fits, with nothing to spare, which is what the outline decision anticipated. \\
P3 & RULED (orchestrator): J1/J2 screw terminal KF128-5.08-2P-AA at 250 V / 24 A ACCEPTED against the requirements' ASSUMED \textgreater{}= 300 V & The \textgreater{}= 300 V figure was an ASSUMED value the requirements-analyst wrote for a part that should never be the limiting element - not an owner requirement and not derived from the board. Actual board voltage never exceeds 30 V, so 250 V is 8x margin; current 24 A vs the \textasciitilde{}2.5 A worst case. No single-piece 5.08 mm 2P THT SCREW terminal at \textgreater{}= 300 V is in stock - the 300 V options change the MECHANISM (two-piece pluggable, or spring-clamp), which is a worse trade than the voltage tier. Requirements section 2 assumption superseded. \\
P3 & THREE part corrections vs architecture/blocks.md, all datasheet-cited; parts.json is the BOM of record for P4 & (1) HF bypass C1 is 220 nF/50 V/X7R per datasheet 9.2.2.6, not blocks.md's 100 nF. (2) C4/C5 case is 1210, not 1206 - no 22 uF/25 V/X7R 1206 part exists with usable stock, and 1210 matches TI's own reference design. (3) {*}{*}C\_VCC 1 uF/16 V was MISSING from blocks.md entirely{*}{*} - datasheet 9.2.2.8 calls it required for proper operation; added as C7. P4 builds from parts.json, so the omission cannot propagate, but blocks.md is stale on these three points and the design doc will show it. \\
P3 & A3 AMENDED by owner: ripple \textless{}= 50 mV pk-pk holds from \textasciitilde{}200 mA to 2 A (CCM region); larger PFM burst ripple accepted below \textasciitilde{}200 mA. +/-3\% DC accuracy is UNCHANGED and still holds at all loads including 0 A. & The LMR33630 DDA package has NO MODE/FPWM pin (8-pin table, none is mode select) and datasheet 8.4.1 states auto-mode only - PFM/diode-emulation at light load, explicitly 'yields larger output voltage ripple'. The only ripple evidence is Fig 8-7, a TYPICAL unguaranteed capture at Vin 12 V / Iout 10 mA, which is neither 0 A nor this board's 18-30 V range: none of the three accept criteria (forced-CCM / MODE pin / datasheet figure at this operating point) is satisfied. Owner ruled to relax rather than override the mode with a preload or re-source. PFM at light load is a deliberate efficiency feature of nearly every modern buck, and on a study article the burst behaviour is itself instructive. P8 sim bounds take the 200 mA threshold. \\
P3 & Min on-time for THIS package is 75-108 ns (DDA), NOT the 68 ns headline - and it is not binding & datasheet\_extract: the Features-bullet 68 ns is the RNX/VQFN value; Table 7.6 gives DDA-specific 75 ns typ / 108 ns max. Required t\_on at the 30 V corner is 5/30/400 kHz = 417 ns, so even the 108 ns max leaves \textasciitilde{}3.9x margin. The P1 selection filter (t\_on\_min \textless{}= 130 ns) is satisfied by the correct number, not only by the optimistic one. \\
P3 & Thermal vias under U1: \textgreater{}= 16 (P1/D4 thermal condition), which also clears the datasheet's own 4x3 = 12 minimum. Figure 10-2's \textasciitilde{}6 vias is NOT the spec. & datasheet\_extract found a genuine intra-datasheet discrepancy: section 10.1.1 text mandates a minimum 4x3 array of 10-mil thermal vias (12) while the worked example Figure 10-2 draws only \textasciitilde{}6. The extraction records 12 as binding text and flags the mismatch so P6/P7 cannot silently copy the smaller figure count. Our own 2L thermal case was computed at \textgreater{}= 16, which is the stricter number and therefore governs. \\
P3 & FOOTPRINT RISK for the librarian/fp\_verify: the datasheet carries TWO DDA mechanical revisions with DIFFERENT exposed-pad sizes - DDA0008J (03/2016) and DDA0008B (09/2025). Extraction used the NEWER DDA0008B (pad copper 2.71 x 3.4 mm). & A pulled or stock footprint built to the 2016 drawing will have the wrong exposed-pad size against the extraction. fp\_verify must be checked against DDA0008B's 1.27 mm pitch, 8 x (0.6 x 1.55 mm) pads, 5.4 mm row spacing and 2.71 x 3.4 mm pad copper. Also note the pinout table itself does not number the pad - 'pin 9' comes only from the mechanical drawing's corner callouts, which is exactly the kind of mismatch that breaks a symbol-to-footprint pin map. \\
P3 & OPEN RISK carried to P4 review: C\_OUT is 2x22uF/25V where TI's own DDA/400 kHz/5 V reference BOM uses 4x22uF/16V. Confirm 2x is inside the part's STABILITY window, not merely its ripple budget. & The part-sourcer justified 2x by finding the 52/72 uF figure in record buck-cmode-internal-comp-cout-window is TI's LOAD-TRANSIENT sizing choice (9.2.2.5), not a stability floor, and bb-buck has no binding transient spec - sound for transient. But the extraction independently reports TI's reference BOM as 4x22uF/16V, and an internally-compensated current-mode loop places its output pole from C\_OUT: 2x22 uF at 25 V derates under 5 V bias to \textasciitilde{}32 uF effective vs \textasciitilde{}66 uF for TI's 4x. Narrow question for the P4 schematic-reviewer to answer FROM THE DATASHEET: is a minimum C\_OUT stated for loop stability (as distinct from transient), and does \textasciitilde{}32 uF effective sit inside it? If not, the fix is cheap now (add C8/C9, same MPN - the 40x30 outline has room) and expensive after P6. \\
P3 & RULED (orchestrator): ACCEPT the vendor U1 footprint as pulled (signal pads 0.6 x 1.3 mm vs DDA0008B nominal 0.6 x 1.55; EP 2.5 x 3.5 vs 2.71 x 3.4). No redraw. & fp\_verify rates it a WARNING, not an error, and pitch 1.27 mm and row spacing 5.4 mm are both exact - the geometry that determines whether the part lands correctly is right. The deviation is small and in the safe direction: shorter toes and a slightly narrower EP mean MORE clearance, not bridging risk; EP area 8.75 vs 9.21 mm2 is a 5\% thermal delta against a case that already carries \textasciitilde{}20\% margin. Decisive argument: this is the LCSC/EasyEDA footprint that JLC's own assembly process is built around, so redrawing to datasheet nominal trades a paper match for a mismatch with the house that actually builds the board. Cost of being wrong is a slightly smaller toe fillet on a bench article, and T8's swap\_new\_fp makes a later change non-destructive. Owner can override. \\
P3 & RULED (orchestrator): ACCEPT the courtyard-cuts-through-pads defect on C0603/R0603/C1210/L1. No redraw. & Documented library-wide easyeda2kicad class (courtyard = body outline only, smaller than the pad span). It is NOT silently breaking placement legality: since the T-fix, pairwise overlap tests the PRECISE shape - declared courtyard UNION pad bbox - so placelib already accounts for the pads the courtyard omits. All 6 footprints HAVE an F.CrtYd (none courtyard-less, which is the case that genuinely degrades the place gate). No clean redraw convention exists yet beyond '0.25 mm beyond pad-field bbox'. Revisit as a library-level fix, not per-board. \\
P3 & No new research opened at the P3 exit despite block:B1 re-firing at the proven floor & The mode's first-call escalation (--maturity-floor proven --research-provisional) re-reports block:B1 as gap because no buck record is yet PROVEN and only bring-up evidence can retire that - it is the same slot already researched at P2 (task block-buck-1, 9 records, 9 verified). Re-opening it would re-research work just done and burn a second per\_run slot for nothing; the mode's own instruction is to re-run WITHOUT the flags after research so the re-run folds in what was learned. The genuinely NEW P3 slot - the part slot (datasheet-layout) - reads COVERED off the C841384 extraction's 16 layout bullets, which is exactly the P3-exit test. Plain re-run: 0 gaps. \\
P4 & ONE schematic agent for the whole board (recipe's small-board amendment), not one per sheet & P2 decision D9 rejected hierarchy: 14 parts / 6 nets on ONE FLAT ROOT SHEET. There is exactly one sheet, so a per-sheet split would spawn one agent and a stitcher for the same file. Recorded per the full-run recipe: 'Small boards may use ONE schematic agent (record it).' \\
P4 & APPROVED after the fact: schematic agent retyped lib/aiee.kicad\_sym U1 pin 7 (BOOT) from power\_in to passive & Evidence-first, and the agent showed its work: with the retype disabled the finished sheet produced EXACTLY ONE ERC finding, power\_pin\_not\_driven on U1 pin 7, and nothing else. /BST has no schematic-visible driver because the boot diode is inside the package. The PWR\_FLAG alternative would clear ERC but drag /BST into netlist\_audit power\_undeclared, and sheets.md s4 deliberately leaves /BST undeclared. Precedent: sbuck-5v3a's AP64350 BST pin carries the same passive typing. Deviations are listed in gen/lib\_pin\_types.py so a later lib\_pull cannot silently re-type wholesale. \\
P4 & E1 CONFIRMED (reviewer, error): C\_OUT 2 x 22 uF is OUTSIDE the datasheet's stability-scoped range. FIX = add C8/C9, same MPN C2918511, to reach 4 x 22 uF. & The earlier phase was right that Eq 6 is a transient bound, but never reached Sec 9.2 p19: 'the internal compensation is optimized for a certain range of external inductance and output capacitance' - quantified ONLY in Table 9-2 p20, whose 400 kHz / 5 V row this board matches on EVERY entry (RFBT 100k, RFBB 24.9k, CBOOT, CVCC, CFF open) except L and C\_OUT, and which specifies 4 x 22 uF RATED. Every 400 kHz row is 4x; 2x appears only at 1400/2100 kHz. TI's own DDA characterization BOM (Table 9-3 p31) used 4 x 22 uF at this exact operating point. Physics: fixed internal Type-II compensation makes fc \textasciitilde{} 1/C\_OUT exactly, so halving C\_OUT roughly doubles crossover toward fSW/2 = 200 kHz, on an all-ceramic bank whose ESR zero sits near 3 MHz and contributes no phase boost. The failure is SILENT at DC and appears only on load/line steps - which is precisely why no gate on this board would have caught it. Derating does not rescue 2x: effective ratio to TI's 16 V bank (\textasciitilde{}0.52) tracks the rated ratio (0.50). Cost \textasciitilde{}\$0.45/board, free now, expensive after P5/P6. \\
P4 & W1 will be FIXED not waived: input bank 2 x 10 uF / 50 V straddles the datasheet's 10 uF EFFECTIVE minimum. Add a third 10 uF (C10, same MPN C2918502). & Sec 9.2.2.6 requires 10 uF MINIMUM and Sec 9.1 p19 defines that as EFFECTIVE, not nameplate. 2 x 10 uF / 50 V gives 10-12 uF effective at 30 V bias, falling to 8.5-10.5 uF with tolerance and the 85-90 C board temperature constraints.json predicts - i.e. BELOW the minimum in the worst corner. It is only a warning because it straddles rather than clearly fails, but the schematic is being regenerated for E1 anyway, the part is \textasciitilde{}\$0.08, and the 40 x 30 mm outline has room. Fixing is strictly better than waiving when the regeneration is already happening. NOTE the mode does NOT argue against this: ultra-bare-bones explicitly INCLUDES the support components the datasheet requires at the stated operating point, so under-populating was the mode violation, not over-populating. \\
P4 & W2 is a LAYOUT carry-forward to P6/P7, not a schematic fix: R2's ground return must be pinned to the AGND/EP reference & Reviewer: nothing currently enforces it - constraints.json's feedback group covers only /FB routing, and the \_review\_enforced item 5 covers only Cin/Cout returns. The exposed pad IS AGND (Table 6-1: all electrical parameters are measured with respect to it), so the divider's bottom-leg return is a reference-accuracy path, not just a ground connection. Carry into the placement and routing spawns explicitly. \\
P4 & CORRECTION to the earlier A3 amendment: DC accuracy at 0 A is NOT unaffected by PFM. Owner ruled BOTH remedies: recentre the divider to 102k/25.5k AND extend the light-load carve-out to cover DC accuracy below \textasciitilde{}200 mA. & sim-analyst finding, from datasheet SNVSAN3F Sec 7.7: VOUT regulation at 5 V is -1.5/+1.5 pct for IOUT \textgreater{}= 1 A but -1.5/{*}{*}+2.5 pct{*}{*} for IOUT 0 A to max - a +1.0 pct light-load adder, stated outright in Sec 8.4.1 (PFM / frequency foldback raises Vout at light load) and visible as a 5.01-5.06 V band in Fig 9-19. A3 explicitly requires 0 A. Stacked on the bench's binding 0-50 C corner: 5.10958 x 1.010 = 5.161 V, MISSING the 5.15 V ceiling by 11 mV. Root cause is NOT the resistors (they own only 0.26 pct of the +2.19 pct budget; VREF and the setpoint offset own the rest) but the +0.32 pct setpoint offset - TI's own Table 9-2 pair sets 5.0161 V, not 5.000. My earlier statement that DC accuracy was unchanged at all loads was wrong; same PFM root cause as the ripple carve-out already accepted. \\
P4 & R1 100k -\textgreater{} 102k (C861068) and R2 24.9k -\textgreater{} 25.5k (C861257): exact 4.000 ratio, 5.000 V nominal. Both E96, both 0.1 pct / 25 ppm, same Yageo RT family (TCR tracking preserved), both verified in stock (2009 / 18845 vs 25 needed). & Recentring moves the stacked worst case from 5.161 V (outside) to \textasciitilde{}5.144 V (inside by 6 mV) and costs nothing - same package, same family, same price class, no library work (R0603 already pulled). It deliberately departs from TI Table 9-2's 100k/24.9k row, which is acceptable because the divider is a SETPOINT item, not a stability item: the ratio sets Vout, and the absolute values matter only for the 1 M CFF threshold (102k is far below it) and for noise, where 102k is the same impedance class. \\
\bottomrule
\end{longtable}
}
\subsection*{Sheet plan}
\subsection*{sheets.md - bb-buck schematic sheet plan, refdes ranges, CANONICAL net names}
\subsubsection*{1. Sheet plan: ONE FLAT ROOT SHEET}
\begin{quote}\small\ttfamily\raggedright
\textbar{} Sheet \textbar{} File \textbar{} Blocks it carries \textbar{} Interface nets at its boundary \textbar{} Refdes ranges \textbar{} `pwr\_base` \textbar{}\\
\textbar{}---\textbar{}---\textbar{}---\textbar{}---\textbar{}---\textbar{}---\textbar{}\\
\textbar{} root (only) \textbar{} `kicad/bb-buck.kicad\_sch` \textbar{} B1 buck converter, J1 input, J2 output, TP1/TP2 probe, H1-H4 mounting \textbar{} none - the board's only interfaces are the two physical terminals (`+VIN`/`GND` at J1, `+5V`/`GND` at J2) \textbar{} `U1-U9`, `L1-L9`, `C1-C19`, `R1-R9`, `J1-J9`, `TP1-TP9`, `H1-H9` \textbar{} {*}{*}1{*}{*} (`\#PWR01..`, `\#FLG01..`) \textbar{}
\end{quote}
\textbf{A hierarchical split is deliberately REJECTED on this board.} Hierarchy exists to keep large schematics readable and to give each sheet its own refdes/\#PWR range; this board has 14 electrical parts and 6 nets on one A4 page. The cost of splitting is concrete and has been paid before: a label inside a child sheet exports as \texttt{/\textless{}sheet\textgreater{}/\textless{}LABEL\textgreater{}}, which silently unhooks every \texttt{constraints.json} entry that spells it \texttt{/\textless{}LABEL\textgreater{}} (the P4 amendment class that cost \texttt{lumina-par}, and the reason \texttt{sbuck-5v3a} also went flat). With one root sheet there are no sheet prefixes anywhere, by construction.

Refdes ranges are therefore trivially unique. They are stated anyway because the rule is "unique ACROSS sheets": {*}{*}if a later amendment ever splits this sheet, the child sheets take \texttt{pwr\_base = 40} and \texttt{80} and keep the prefix ranges above{*}{*}, and every net name in s2 must be re-checked against the new export.

\subsubsection*{2. CANONICAL NET NAMES - these are now binding}
\texttt{research/power.json} proposed \texttt{VIN / SW / +5V / GND}. {*}{*}P2 makes them canonical here, and two of the four change spelling.{*}{*} These are the names the schematic MUST produce and the names \texttt{constraints.json} uses.

\begin{quote}\small\ttfamily\raggedright
\textbar{} Canonical net \textbar{} KiCad mechanism \textbar{} Spans \textbar{} Declared in constraints.json as \textbar{}\\
\textbar{}---\textbar{}---\textbar{}---\textbar{}---\textbar{}\\
\textbar{} {*}{*}`+VIN`{*}{*} \textbar{} power symbol (VALUE = net name) -\textgreater{} exports BARE \textbar{} J1.1 -\textgreater{} C1/C2/C3 -\textgreater{} U1 VIN pin \textbar{} `power` 1.1 A, `voltages` 30 V \textbar{}\\
\textbar{} {*}{*}`GND`{*}{*} \textbar{} power symbol -\textgreater{} exports BARE \textbar{} everything; B.Cu pour \textbar{} `power` 2.6 A `plane\_fed`, `planes` F.Cu + B.Cu, `thermal` heatsink net \textbar{}\\
\textbar{} {*}{*}`+5V`{*}{*} \textbar{} power symbol -\textgreater{} exports BARE \textbar{} L1 -\textgreater{} C4/C5 -\textgreater{} R1 -\textgreater{} J2.1 \textbar{} `power` 2.6 A \textbar{}\\
\textbar{} {*}{*}`/SW`{*}{*} \textbar{} root-sheet LOCAL LABEL -\textgreater{} exports with ONE leading slash \textbar{} U1 SW pin -\textgreater{} L1 -\textgreater{} TP1 \textbar{} `power` 2.6 A `pdn:false`, `high\_speed`, `voltages` 30 V \textbar{}\\
\textbar{} {*}{*}`/FB`{*}{*} \textbar{} root-sheet local label \textbar{} R1/R2 midpoint -\textgreater{} U1 FB pin \textbar{} {*}{*}deliberately NOT declared{*}{*} (see s4) \textbar{}\\
\textbar{} {*}{*}`/BST`{*}{*} \textbar{} root-sheet local label \textbar{} U1 BOOT pin -\textgreater{} C6 -\textgreater{} `/SW` \textbar{} {*}{*}deliberately NOT declared{*}{*} (see s4) \textbar{}
\end{quote}
Why \texttt{+VIN} and not \texttt{VIN}: it is a rail feeding a power pin and decoupled by the input bank, so it is a power symbol, and a power symbol's exported name is bare. Why \texttt{/SW} and not \texttt{SW}: it is a root-sheet local label, and KiCad prefixes those with a slash on export. {*}{*}A power symbol WINS over a coincident label{*}{*}, so P4 must not mix the two mechanisms on one node.

\textbf{Enforcement, not hope:} at P4 run \texttt{netlist\_audit.py --net kicad/bb-buck.net --constraints architecture/constraints.json}. A net spelled differently in the schematic surfaces as \texttt{missing\_net} (error). Without that run, a misspelling is silent: \texttt{check\_current} finds no copper on a net that does not exist and reports nothing.

\subsubsection*{3. Contents of the root sheet}
\begin{quote}\small\ttfamily\raggedright
\textbar{} Refdes \textbar{} Part \textbar{} Net attachments \textbar{}\\
\textbar{}---\textbar{}---\textbar{}---\textbar{}\\
\textbar{} `U1` \textbar{} 36 V-class synchronous integrated-FET buck, 3 A-class, exposed pad, 400 kHz, internally compensated \textbar{} VIN=`+VIN`, GND/PGND/EP=`GND`, SW=`/SW`, BOOT=`/BST`, FB=`/FB`, EN per datasheet (see below) \textbar{}\\
\textbar{} `C1` \textbar{} 100 nF 50 V X7R - HF input bypass, {*}{*}`role: reg\_input`{*}{*} \textbar{} `+VIN` / `GND` \textbar{}\\
\textbar{} `C2` `C3` \textbar{} 10 uF 1210 50 V X7R - input bank \textbar{} `+VIN` / `GND` \textbar{}\\
\textbar{} `C6` \textbar{} 100 nF - bootstrap (value/rating per datasheet) \textbar{} `/BST` / `/SW` \textbar{}\\
\textbar{} `L1` \textbar{} 15 uH shielded/composite, \textless{}= 40 mOhm, Isat \textgreater{}= \textasciitilde{}6.6 A \textbar{} `/SW` / `+5V` \textbar{}\\
\textbar{} `C4` `C5` \textbar{} 22 uF 1206 25 V X7R - output bank \textbar{} `+5V` / `GND` \textbar{}\\
\textbar{} `R1` `R2` \textbar{} FB divider, 0.1 \% / 25 ppm (`R1` = top, `+5V`-\textgreater{}`/FB`; `R2` = bottom, `/FB`-\textgreater{}`GND`) \textbar{} as listed \textbar{}\\
\textbar{} `J1` \textbar{} 2-pole 5.08 mm screw terminal, THT \textbar{} 1=`+VIN`, 2=`GND` \textbar{}\\
\textbar{} `J2` \textbar{} same part as J1 \textbar{} 1=`+5V`, 2=`GND` \textbar{}\\
\textbar{} `TP1` \textbar{} SW probe pad (A4) - INSIDE the SW copper, not on a spur \textbar{} `/SW` \textbar{}\\
\textbar{} `TP2` \textbar{} GND probe pad, adjacent to TP1 \textbar{} `GND` \textbar{}\\
\textbar{} `H1`-`H4` \textbar{} M3 clearance holes, 3.2 mm \textbar{} mechanical (no net) \textbar{}
\end{quote}
\textbf{EN: read the part's own behaviour before adding anything.} Most parts in this class auto-start from an internal pull-up, so EN ties to \texttt{+VIN} or floats and NO parts are added. No UVLO divider - mode-excluded, and a bench supply has no cable-droop motorboating case (0.62 A over 2 m of lead is \textasciitilde{}0.1 V). No soft-start components: internal soft-start makes Cout inrush \texttt{C x Vout/tss} \textasciitilde{} 65 mA, and A2 means the supply ramps from zero anyway.

If the FB pin has no divider because P3 found a FIXED 5 V part, \texttt{R1}/\texttt{R2} and \texttt{/FB} disappear and the sense point ties straight to the output - {*}{*}that is the preferred outcome and it removes the whole s5 tolerance budget{*}{*}, but do not synthesise it from an adjustable part's datasheet figure.

\subsubsection*{4. Two nets deliberately left OUT of \texttt{constraints.json}}
\texttt{/FB} and \texttt{/BST} are declared nowhere. This is intentional:

\begin{itemize}
\item \texttt{rules\_gen} buckets every \texttt{power} entry into a netclass {*}{*}at that net's own
IPC-2152 width{*}{*} (this used to flatten every power net to the widest width;
the current script buckets per net - verified in \texttt{rules\_gen.net\_classes}).
Even so, a \texttt{power} entry on \texttt{/FB} would give the feedback node a
minimum-width rule and pull it into the \texttt{check\_pdn} decoupling inventory,
and \texttt{/FB} is a high-impedance sense node that wants to be SHORT and THIN
and nowhere near \texttt{/SW}.
\item \texttt{/BST} carries only gate charge. Nothing to size, nothing to decouple.
\end{itemize}
The two rules that DO apply to them are layout rules, carried in s5 below and in \texttt{constraints.json} \texttt{placement.groups} / \texttt{separation}, not width rules.

\subsubsection*{5. Placement groups this sheet implies (they become P6 groups)}
\begin{itemize}
\item \textbf{`hotloop`} (anchor \texttt{U1}; \texttt{C1}, \texttt{C2}, \texttt{C3}, \texttt{C6}): C1 nearest the VIN
pin, then the bulk ceramics, then U1 VIN -\textgreater{} U1 PGND, on F.Cu, {*}{*}smallest
enclosed area achievable{*}{*}. This is the loop that spends \textasciitilde{}37 mV of the
50 mV ripple budget. It must be placed by explicit \texttt{place\_edit} and LOCKED
before \texttt{place\_anneal} runs - an annealer optimises a cost function and
cannot discover a hot loop.
\item \textbf{`output`} (anchor \texttt{L1}; \texttt{C4}, \texttt{C5}): Cout ground and Cin ground return
to the GND pour at SEPARATE points near their own pins - never a shared
narrow neck (\texttt{buck-cin-co-ground-separation}).
\item \textbf{`feedback`} (anchor \texttt{R1}; \texttt{R2}): at the FB pin, sense point AFTER the
output caps, \textgreater{}= 5 mm from L1 and off the \texttt{/SW} side of the package.
\item \textbf{`probe`} (anchor \texttt{TP1}; \texttt{TP2}): TP1 inside the existing \texttt{/SW} pour;
TP2's ground returns straight down to the B.Cu pour beneath it.
\end{itemize}
\texttt{/SW} itself: a short WIDE pour, 1.6-2.0 mm wide, \textless{}= 8 mm long, \textless{}= 40 mm\textasciicircum{}2 total, F.Cu only, never near a board edge or a screw terminal. IPC-2152 constrains the WIDTH (a floor); EMI constrains the AREA (a ceiling that nothing in the toolchain enforces). TP1's pad counts against that 40 mm\textasciicircum{}2.

Full architecture narrative (not inlined; see the artifact index): \texttt{architecture/blocks.md}, \texttt{architecture/decisions.md}, \texttt{architecture/power\_tree.md}, \texttt{architecture/stackup.md}.
\section{Schematic}
Gate \texttt{erc}: \emph{not recorded}.
The full schematic PDF follows.
\includepdf[pages=-]{reports/schematic.pdf}
\section{Layout}
\emph{Pending --- produced at P6.}
\section{Verification}
\emph{Pending --- produced at P8.}
\section{DFM and Fabrication}
\emph{Pending --- produced at P9.}
\section{Run Record (Appendix)}
\subsection*{Phase digests}
\paragraph*{P0}
\begin{quote}\small\ttfamily\raggedright
\# P0 Intake - digest (2026-08-15)\\
~\\
- bb-buck: non-isolated buck, 18-30 V DC in (24 nom) -\textgreater{} 5 V @ 0-2 A out, screw\\
  terminal both ends, bench supply + resistive load.\\
- Mode {*}{*}ultra-bare-bones{*}{*} recorded as a decision: buck block + datasheet-required\\
  support only. Defaults applied (qty 5, JLC PCBA single-sided, 0-50 C bench, no\\
  enclosure, fewest honest layers, smallest honest outline, stop at P9).\\
- requirements.md by `requirements-analyst`; `check\_requirements.py` PASS 9/9, 0\\
  violations, before and after fold-in.\\
- Section 8 closed with NO safety question open, derived: bench source (no battery),\\
  resistive load (no motor), \textasciitilde{}2.5 A peak (\textless{} 3 A), 30 V is not \textgreater{}30 V and \textless{} 60 V ES1.\\
- 4 questions asked in ONE batch, owner took all four defaults -\textgreater{} now requirements:\\
  A1 30 V hard max operating, part abs-max \textgreater{}= 36 V (headroom in the PART, no clamp);\\
  A2 no live hot-plug (drops the \textasciitilde{}60 V ring case); A3 +/-3 \% DC over full line+load,\\
  \textless{}= 50 mV pk-pk; A4 one switch-node probe pad + adjacent GND pad, nothing else.\\
- Left open for P1/P2: sync vs async rectification; 2L vs 4L (2L assumed, revisable\\
  on thermal grounds at \textasciitilde{}0.8-1.8 W / 50 C ambient).\\
- Fable tier out of credits -\textgreater{} fable/high spawns run opus/xhigh (ledgered).
\end{quote}
\paragraph*{P1}
\begin{quote}\small\ttfamily\raggedright
\# P1 Research - digest (2026-08-15)\\
~\\
- Roster: power-architect + component-scout(buck-regulator) only. No interface-spec\\
  (screw terminals bind no standard), no reference-design (P3 part-level coverage\\
  research produces the vendor layout digest for the CHOSEN IC instead).\\
- Topology ruled {*}{*}synchronous integrated-FET, \textasciitilde{}500 kHz (350-600), L = 10 uH{*}{*}. Async\\
  rejected: +0.42 W (+33 \%) plus a 0.77 W diode at Tj \textasciitilde{}126 C vs a 125 C rating (D=0.17\\
  means it conducts 83 \% of every cycle at 30 V). \textgreater{}= 1 MHz rejected: 1.58-1.93 W and\\
  t\_on(30 V) collides with min-on-time.\\
- Part envelope for P3: Vin abs-max \textgreater{}= 36 V, {*}{*}Rds\_LS \textless{}= 85 mOhm (rank on LS not HS+LS;\\
  at D=0.167 the LS FET conducts 5x longer){*}{*}, t\_on\_min \textless{}= 130 ns, Tj 150 C, exposed pad\\
  mandatory, P\_U1 \textless{}= 0.95 W at 30 V/2 A. A 3 A-class part, not a marginal 2 A one.\\
- {*}{*}2L CONDITIONAL{*}{*} - the screen passes by 2 C (0.92 W x 73.8 C/W = 67.9 vs dt\_c 70).\\
  Needs sync + Rds\_LS \textless{}= 85 mOhm + outline \textgreater{}= \textasciitilde{}1000 mm\textasciicircum{}2 + \textgreater{}= 16 thermal vias + Tj 150 C.\\
  Layer count FOLLOWS THE PART; re-run at P3 before P5 fixes the stackup. Mode note:\\
  "smallest honest outline" has a numeric floor - the outline IS the radiator.\\
- Scout, 5 candidates, every figure datasheet-sourced: top {*}{*}LMR33630ADDAR{*}{*} C841384\\
  (sync, 38 V, 400 kHz, \$0.74@5, stock 6885); AOZ1284PI C48060 async 40 V \$0.85;\\
  TPS54560BDDAR C1850354 65 V \$1.06. AP63357Q out at 35 V abs-max (1 V short of A1);\\
  LMR33630's 1.4/2.1 MHz siblings pulse-skip at 30 V - only the 400 kHz "A" clears it;\\
  XL1509 is the only Basic part but its own +/-4 \% already breaks A3's +/-3 \%.\\
- Convergence: the scout's top pick sits inside the architect's envelope.
\end{quote}
\paragraph*{P2}
\begin{quote}\small\ttfamily\raggedright
\# P2 Architecture - digest (2026-08-15)\\
~\\
- One block (B1 sync buck) + J1/J2 screw terminals + TP1/TP2 SW probe pair + 4x M3.\\
  14 parts, 6 nets, ONE FLAT ROOT SHEET. Canonical nets: +VIN /SW +5V GND.\\
  Stackup JLC2313\_1.6 2L 1 oz CONDITIONAL (trigger: P\_U1 \textgreater{} 0.95 W = Rds\_LS \textgreater{} \textasciitilde{}110\\
  mOhm at 400 kHz). {*}{*}Outline 40 x 30 mm FINAL{*}{*} - owner-confirmed at H1.\\
- fsw 400 kHz / L 15 uH. P1's fixed-output preference LOST to stock reality -\textgreater{}\\
  adjustable + 0.1 \%/25 ppm divider, closing A3 on worst-case SUM. lint 0/0.\\
- Architect corrected 3 briefing premises, all verified (rules\_gen already buckets\\
  per IPC-2152 width; planes{[}{]} REPLACES defaults and rejects `\_note`; pdn:false also\\
  skips check\_irdrop) -\textgreater{} LEARNINGS.md entry, incl. a latent sbuck-5v3a defect.\\
- Coverage at the mode floor: 1 slot / 10 classes all gap (8 maturity-only, 3\\
  substantive). Mapper skipped as a provable no-op. Research block-buck-1: 3 TI\\
  sources, 12+ pages read visually, 9 records. Reader REFUTED 1 (record said 0.3 uH,\\
  figure prints {*}{*}9.3 uH{*}{*} - 31x, mis-sized the damping criterion); corrected,\\
  re-read fresh -\textgreater{} {*}{*}9 verified / 0 refuted{*}{*}. Re-run: 9 covered / 1 provisional\\
  (inrush) / 0 gap.\\
- Design-relevant: hot-plug at 30 V rings to 45-49 V vs 38 V abs-max (A2 is what\\
  makes the clamp-free mode safe); the EP is {*}{*}AGND{*}{*}, so the via array is a\\
  loop-accuracy item; 2L is a documented deviation from the datasheet's layout rule 5\\
  at \textasciitilde{}5 dBuV/m -\textgreater{} P6/P7 must keep the bottom pour unbroken under hot loop, SW and L.\\
- H1 APPROVED. Promotion of the 9 records deferred to run close with bring-up.
\end{quote}
\paragraph*{P3}
\begin{quote}\small\ttfamily\raggedright
\# P3 Parts + Library - digest (2026-08-15)\\
~\\
- parts.json: 10 distinct parts / 13 refdes, \textasciitilde{}\$1.63/board, 1 Basic + 9 Extended.\\
  U1 has NO drop-in alternate (async candidates D1-rejected; frequency siblings\\
  cascade into L/C\_OUT re-derivation).\\
- {*}{*}D4 MET - 2 LAYERS CONFIRMED{*}{*}: real Rds\_LS 66 typ / 110 max mOhm over full temp\\
  -\textgreater{} P\_U1 0.57 typ / 0.76 W worst-case vs the 0.95 W gate, \textasciitilde{}20 \% margin (the\\
  placeholder model predicted 2 C). D3 met (15 uH + 2x22 uF inside both windows).\\
  D11 met (Isat 8 A vs 6.6 A floor, DCR 27 mOhm, 12.3 mm body).\\
- {*}{*}A3 AMENDED by owner{*}{*}: \textless{}= 50 mV holds 200 mA - 2 A, PFM burst ripple accepted\\
  below. DDA has no MODE pin, auto-mode only (8.4.1). DC accuracy unchanged.\\
- Extraction validated (9 pins, EP = {*}{*}AGND{*}{*}). Traps: the 68 ns min-on-time headline\\
  is the VQFN part's (DDA 75-108 ns vs our 417 ns need); 10.1.1's 4x3=12 vias\\
  contradicts Fig 10-2's \textasciitilde{}6 - our \textgreater{}= 16 governs. Two DDA revs with different EP sizes.\\
- Library 10/10 pulled, 6/6 load-check, 27 pins retyped, no B.Cu lands, no netless PTH\\
  in the EP (LM5017 precedent did not repeat). RULED: accept the vendor U1 land\\
  (warning-level; it is what JLC's process is built around) and the library-wide\\
  courtyard-cuts-pads class (placelib unions pad bbox).\\
- {*}{*}Coverage P3: 1 covered (part slot) / 1 provisional (inrush) / 0 gaps.{*}{*}\\
- OPEN RISK to P4 review: C\_OUT 2x22 uF vs TI's 4x22 uF reference - stability or only\\
  ripple?
\end{quote}
\paragraph*{P4}
\begin{quote}\small\ttfamily\raggedright
\# P4 Schematic - digest (2026-08-15)\\
~\\
- ONE schematic agent (flat root sheet, recorded deviation). Generator\\
  `kicad/gen/root.py` is the source; the .kicad\_sch is build output.\\
- {*}{*}Gates: erc 0 errors / 0 warnings; netlist\_audit exit 0; sim PASS\\
  (0 err / 0 warn, 10 bounds, 18 measures).{*}{*} 22 components, 8 nets.\\
- Cleared by the reviewer on datasheet evidence: EN-\textgreater{}+VIN direct (Table 6-1\\
  sanctions it, and EN abs-max is VIN+0.3 so it cannot be violated by\\
  construction), PG open no-connect, EP=AGND tied to GND (the datasheet\\
  REQUIRES the tie), CFF correctly open (confirmed three ways), full abs-max\\
  sweep at 30 V, ton 417 ns vs 108 ns, Isat 8 A vs the strict 5.05 A criterion.\\
- {*}{*}Fix pass, 3 changes.{*}{*} E1 (error): C\_OUT 2 -\textgreater{} {*}{*}4 x 22 uF{*}{*} - 2x is outside\\
  the internal compensation's optimised range (Sec 9.2 p19, quantified only in\\
  Table 9-2 p20 whose 400 kHz/5 V row is 4x; every 400 kHz row is 4x). Failure\\
  would have been SILENT at DC, visible only on load/line steps. W1: input bulk\\
  2 -\textgreater{} {*}{*}3 x 10 uF{*}{*} (the 10 uF minimum is EFFECTIVE, and 2x fell to \textasciitilde{}8.5 uF in\\
  the worst corner). A3: divider {*}{*}100k/24.9k -\textgreater{} 102k/25.5k{*}{*}, exact 4.000\\
  ratio, 5.000 V nominal. +\$0.52/board.\\
- {*}{*}A3 amended a second time and requirements.md updated.{*}{*} Sec 7.7 gives\\
  -1.5/+2.5 \% regulation for IOUT 0 A to max (vs +/-1.5 \% above 1 A) - the same\\
  PFM root cause as the ripple carve-out. Stacked worst case was 5.161 V vs the\\
  5.15 V ceiling; recentring moved it to 5.144 V. My earlier "DC accuracy\\
  unaffected" statement was wrong and is retracted in requirements.md.\\
- Sim calibration proved the light-load warning bound is load-bearing: a revert\\
  to 100k/24.9k produces ZERO errors and is caught ONLY by `vout\_ll0a`. All four\\
  single-E96-step mutations trip at error severity.\\
- constraints.json placement.groups repaired by the orchestrator (hotloop was\\
  missing C10, output was missing C8/C9 and still said "1206"); W2 carried into\\
  the feedback group note - R2's return must pin to the AGND/EP reference.\\
- Library: C861068 has NO EasyEDA CAD record at all (lib\_pull reports that\\
  identically to a rate-limit) - symbol cloned from its same-family sibling.
\end{quote}
\subsection*{Gate attempt history}
\emph{(no entries)}
\subsection*{Issues}
\emph{(no entries)}
\subsection*{Human checkpoints}
{\small
\begin{longtable}{lllp{2.6cm}p{6.4cm}}
\toprule
\textbf{Id} & \textbf{Phase} & \textbf{Status} & \textbf{When} & \textbf{Note} \\
\midrule
\endhead
1 & P2 & approved & 2026-08-15T18:56:25 & Architecture approved. Outline 40x30 mm CONFIRMED by owner as final (irreversible at P5 board\_init). Record promotion deferred to run close with bring-up evidence - verified workspace records still feed P3/P6/P7 spawns meanwhile; inrush stays provisional until then. \\
2 & P4 & not reached &  &  \\
3 & P6 & not reached &  &  \\
4 & P8 & not reached &  &  \\
5 & P10 & not reached &  &  \\
\bottomrule
\end{longtable}
}
\section{Artifact Index}
{\small
\begin{longtable}{p{9cm}rp{4cm}}
\toprule
\textbf{File} & \textbf{Size} & \textbf{Note} \\
\midrule
\endhead
\texttt{state.json} & 42,099 B &  \\
\texttt{kicad/bb-buck.kicad\_sch} & 82,705 B &  \\
\texttt{brief/brief.md} & 605 B &  \\
\texttt{requirements.md} & 18,721 B &  \\
\texttt{architecture/blocks.md} & 10,091 B &  \\
\texttt{architecture/constraints.json} & 18,896 B &  \\
\texttt{architecture/decisions.md} & 9,227 B &  \\
\texttt{architecture/power\_tree.md} & 10,485 B &  \\
\texttt{architecture/sheets.md} & 7,162 B &  \\
\texttt{architecture/stackup.md} & 7,324 B &  \\
\texttt{reports/gate-erc.json} & 711 B &  \\
\texttt{reports/gate-sim.json} & 675 B &  \\
\texttt{reports/review-schematic.json} & 2,943 B &  \\
\texttt{reports/review-schematic.md} & 24,242 B &  \\
\texttt{reports/schematic.pdf} & 76,058 B &  \\
\texttt{reports/sim.json} & 1,214 B &  \\
\texttt{reports/top.net} & 34,564 B &  \\
\bottomrule
\end{longtable}
}
\end{document}
